% !TEX root = ../DoAn.tex
\documentclass[../DoAn.tex]{subfiles}
\begin{document}

Chương này trình bày quá trình hiện thực hóa thiết kế đã phân tích ở Chương~\ref{chapter:Experiment}.
Do báo cáo được giới hạn trong năm chương, nội dung cài đặt, triển khai, kiểm thử và đánh giá
được tổ chức chung trong một chương. Cách tổ chức này giúp thể hiện mạch đầy đủ từ thiết kế
đến hệ thống chạy được, đồng thời vẫn tách rõ phần hiện thực và phần kiểm chứng bằng các mục
riêng biệt.

% =======================================================================
\section{Môi trường phát triển và cấu trúc dự án}
\label{section:5.1}

Hệ thống được phát triển theo mô hình monorepo để thuận tiện cho việc quản lý dependency,
build, kiểm thử và chạy toàn bộ stack trong phạm vi đồ án. Backend sử dụng Java 21, Spring
Boot 3.5.13, Spring Cloud 2025.0.0 và Maven Wrapper. Frontend sử dụng Next.js App Router,
TypeScript và Server Actions. Toàn bộ môi trường runtime được chuẩn hóa bằng Docker Compose.

\begin{table}[H]
\centering
\caption{Môi trường và công cụ phát triển}
\label{table:dev-env}
\begin{compacttable}
\begin{tabularx}{\textwidth}{|L{3cm}|L{4.1cm}|Y|}
\hline
\textbf{Nhóm} & \textbf{Công cụ/Phiên bản} & \textbf{Vai trò trong đồ án} \\ \hline
Ngôn ngữ backend & Java 21 & Phát triển 13 business service và 3 infrastructure service \\ \hline
Framework backend & Spring Boot 3.5.13, Spring Cloud 2025.0.0 & REST API, Gateway, Eureka, Config, OpenFeign, Kafka, Security, Actuator \\ \hline
Build backend & Maven Wrapper & Build multi-module, quản lý dependency và plugin tập trung \\ \hline
Frontend & Next.js App Router, TypeScript & Storefront, Admin Panel, BFF/proxy, Server Actions \\ \hline
Container & Docker Compose & Khởi động local stack và production-like stack trên EC2 \\ \hline
Kiểm thử & JUnit 5, Testcontainers, k6, Bash script & Unit, integration, smoke, performance và resilience test \\ \hline
Triển khai & AWS EC2, Caddy, GitHub Actions, GHCR & Production-like single-host deployment và build/push image \\ \hline
\end{tabularx}
\end{compacttable}
\end{table}

Về cấu trúc mã nguồn, repository gồm module \path|common|, 13 business service
như \path|identity|, \path|user|, \path|product|, \path|inventory|, \path|cart|,
\path|order|, \path|payment|, \path|voucher|, \path|notification|, \path|review|,
\path|search|, \path|content| và \path|flash-sale|; ba service hạ tầng
\path|api-gateway|, \path|discovery-server|, \path|config-server|; và thư mục
\path|frontend|. Mỗi service backend có Dockerfile, cấu hình riêng và nếu có trạng thái
thì có Flyway migration quản lý schema.

% \begin{figure}[H]
%     \centering
%     \includegraphics[width=0.85\textwidth]{DUONG_DAN_ANH}
%     \caption{Cấu trúc thư mục dự án và các module chính}
%     \label{fig:project-structure}
% \end{figure}

% =======================================================================
\section{Cài đặt các thành phần backend cốt lõi}
\label{section:5.2}

\subsection{Spring Cloud layer}

Ba service hạ tầng tạo thành lớp điều phối cho toàn bộ backend. \texttt{api-gateway} là điểm
vào duy nhất cho client, chịu trách nhiệm định tuyến, xác thực JWT, gắn trusted headers và
rate limiting. \texttt{discovery-server} sử dụng Eureka để các service đăng ký và phát hiện
lẫn nhau theo tên logic thay vì hard-code địa chỉ. \texttt{config-server} cung cấp cấu hình
tập trung từ classpath, giúp các service lấy cấu hình theo profile khi khởi động.

% \begin{figure}[H]
%     \centering
%     \includegraphics[width=0.85\textwidth]{DUONG_DAN_ANH}
%     \caption{Spring Cloud layer: Gateway, Discovery Server và Config Server}
%     \label{fig:spring-cloud-layer}
% \end{figure}

\subsection{Database per Service và Flyway}

Các service stateful sở hữu database logic riêng trên PostgreSQL. Cách tổ chức này đảm bảo
nguyên tắc Database per Service: service khác không đọc/ghi trực tiếp vào database không
thuộc quyền sở hữu của mình. Schema được quản lý bằng Flyway migration và Hibernate chạy với
chế độ validate, nhờ đó sai lệch giữa entity và schema được phát hiện sớm khi service khởi
động.

\begin{table}[H]
\centering
\caption{Map service với lớp lưu trữ chính}
\label{table:service-storage-map}
\begin{compacttable}
\begin{tabularx}{\textwidth}{|L{3.1cm}|L{3cm}|Y|}
\hline
\textbf{Service} & \textbf{Lưu trữ chính} & \textbf{Ghi chú} \\ \hline
identity-service & Keycloak & Tạo user trong realm, phát sự kiện đăng ký \\ \hline
user-service & PostgreSQL & Hồ sơ người dùng, địa chỉ, dữ liệu admin user \\ \hline
product-service & PostgreSQL, Redis & Catalog, category, brand, cache product list/detail \\ \hline
inventory-service & PostgreSQL & Tồn kho, reserved quantity, stock movement \\ \hline
cart-service & Redis & Giỏ hàng user/guest với TTL \\ \hline
order-service & PostgreSQL, Kafka & Order, order item, outbox, processed events \\ \hline
payment-service & PostgreSQL, Kafka & Payment, payment outbox, processed events \\ \hline
search-service & Elasticsearch & Read model phục vụ full-text search \\ \hline
flash-sale-service & PostgreSQL, Redis, Kafka & Campaign, Redis Lua reservation, reconciliation \\ \hline
\end{tabularx}
\end{compacttable}
\end{table}

\subsection{Saga Orchestration, Outbox và Idempotent Consumer}

Luồng đặt hàng là phần thể hiện rõ nhất các mẫu thiết kế phân tán. \texttt{order-service}
đóng vai trò orchestrator: tạo đơn, đặt chỗ voucher, phát sự kiện \texttt{order-created},
nhận kết quả từ inventory/payment và quyết định xác nhận hoặc hủy đơn. Để tránh lỗi dual-write,
\texttt{order-service} ghi sự kiện vào bảng \texttt{outbox} trong cùng transaction với dữ liệu
đơn hàng. \texttt{OutboxPoller} quét 1 giây/lần, lấy batch chưa xử lý bằng
\texttt{FOR UPDATE SKIP LOCKED}, gửi Kafka và đánh dấu \texttt{processed=true}.

% \begin{figure}[H]
%     \centering
%     \includegraphics[width=0.85\textwidth]{DUONG_DAN_ANH}
%     \caption{Snippet hoặc sơ đồ OutboxPoller đọc bảng outbox và phát Kafka}
%     \label{fig:outbox-poller}
% \end{figure}

Ở phía consumer, các service như order, inventory và payment sử dụng bảng
\texttt{processed\_events} để lưu \texttt{eventId} đã xử lý. Nếu Kafka gửi lại cùng một event,
consumer bỏ qua thay vì thực thi nghiệp vụ lần hai. Cách làm này chấp nhận at-least-once
delivery ở tầng message broker nhưng đạt hiệu ứng exactly-once processing ở tầng ứng dụng.

\subsection{Thanh toán VNPAY}

\texttt{payment-service} tích hợp VNPAY sandbox theo ba bước chính. Khi người dùng chọn thanh
toán online, service tạo payment ở trạng thái \texttt{PENDING}, sinh URL thanh toán và ký tham
số bằng HMAC-SHA512. Khi VNPAY gọi IPN hoặc người dùng quay lại Return URL, service xác minh
chữ ký, kiểm tra số tiền, cập nhật trạng thái payment và ghi sự kiện vào
\path|payment_outbox|. \path|PaymentEventOutboxPoller| phát sự kiện
\path|payment-success| hoặc \path|payment-failed| để \path|order-service| hoàn tất Saga.
Nếu payment quá hạn, \path|PaymentTimeoutScheduler| chuyển trạng thái sang \path|TIMEOUT|.

\subsection{Flash sale high-concurrency}

\path|flash-sale-service| xử lý mua flash sale bằng Redis Lua script để thao tác kiểm tra
buyer set, kiểm tra stock và giảm slot trong một lệnh nguyên tử. Redis key chính có dạng
\path|flash-sale:{campaignId}:stock| và \path|flash-sale:{campaignId}:buyers|, được
đặt TTL theo thời điểm kết thúc chiến dịch. Sau khi giữ slot thành công, service phát event
\path|flash-sale-order-requested|; nếu tạo đơn thất bại, compensation script hoàn lại slot
và xóa buyer khỏi set.

% \begin{figure}[H]
%     \centering
%     \includegraphics[width=0.85\textwidth]{DUONG_DAN_ANH}
%     \caption{Redis Lua reservation và compensation cho flash sale}
%     \label{fig:redis-lua-code}
% \end{figure}

Hai scheduler hỗ trợ flash-sale vận hành ổn định. \path|CampaignScheduler| chạy 5 giây/lần,
dùng Redis distributed lock để chỉ một instance kích hoạt chiến dịch đến hạn, khôi phục Redis
stock key bị mất và kết thúc campaign hết hạn. \path|ReconciliationScheduler| chạy 300 giây/lần,
đối chiếu \path|soldCount| trong flash-sale-service với số order flash-sale thực tế ở
order-service, từ đó phát hiện và bù slot mồ côi khi cần.

% =======================================================================
\section{Cài đặt bảo mật và phân quyền}
\label{section:5.3}

Bảo mật được thực thi tập trung tại API Gateway. Client gửi access token dạng
\path|Authorization: Bearer <token>| tới Gateway; Gateway xác minh chữ ký RS256 qua JWKS
của Keycloak, kiểm tra thời hạn token, trích xuất userId/email/roles rồi gắn vào ba trusted
headers: \path|X-User-Id|, \path|X-User-Email|, \path|X-User-Roles|. Backend service
chỉ tin headers do Gateway đưa vào trong mạng nội bộ Docker, không mở cổng trực tiếp ra ngoài.

\begin{table}[H]
\centering
\caption{Chính sách phân quyền tại API Gateway}
\label{table:gateway-auth-policy}
\begin{compacttable}
\begin{tabularx}{\textwidth}{|L{4.2cm}|L{2.8cm}|Y|}
\hline
\textbf{Nhóm endpoint} & \textbf{Quyền yêu cầu} & \textbf{Ghi chú} \\ \hline
\path|GET /api/products/**|, \path|GET /api/search/**|, \path|GET /api/content/**| & Public & Cho phép người dùng duyệt catalog, search và nội dung CMS \\ \hline
\path|/api/cart/**| & Public hoặc token & Hỗ trợ guest cart bằng session id và user cart bằng token \\ \hline
\path|POST /api/orders|, user profile, review create & Authenticated & Yêu cầu token hợp lệ và trusted headers \\ \hline
\path|/api/orders/admin/**|, \path|/api/users/admin/**|, \path|/api/inventory/admin/**|, \path|/api/reviews/admin/**| & \path|ROLE_ADMIN| & Endpoint quản trị bổ sung trong Phase 14 \\ \hline
Product/voucher/flash-sale/content mutations & \path|ROLE_ADMIN| & POST/PUT/DELETE được bảo vệ theo method và path \\ \hline
\path|POST /api/payments/vnpay/ipn| & Public có HMAC & Public vì do VNPAY gọi, nhưng bắt buộc xác minh chữ ký \\ \hline
\end{tabularx}
\end{compacttable}
\end{table}

Ngoài xác thực và RBAC, Gateway áp dụng Redis token bucket rate limiter cho các endpoint nhạy
cảm. Luồng đặt hàng được giới hạn 10 request/giây, burst 20; luồng mua flash sale được giới
hạn 3 request/giây, burst 5. Rate limiter không thay thế kiểm soát nghiệp vụ nhưng giúp giảm
tải và hạn chế automation trong các luồng có rủi ro cao.

% \begin{figure}[H]
%     \centering
%     \includegraphics[width=0.85\textwidth]{DUONG_DAN_ANH}
%     \caption{Security flow: Client, Gateway, Keycloak và backend service}
%     \label{fig:security-flow-ch5}
% \end{figure}

% =======================================================================
\section{Cài đặt frontend Storefront và Admin Panel}
\label{section:5.4}

Frontend được xây dựng trong thư mục \path|frontend/| bằng Next.js App Router. Ứng dụng có
hai khu vực chính: storefront cho khách hàng và Admin Panel cho quản trị viên. Storefront hỗ
trợ duyệt sản phẩm, xem chi tiết, tìm kiếm, giỏ hàng, checkout, xem đơn hàng, thanh toán VNPAY
và các màn hình flash-sale, content, review phục vụ demo. Admin Panel nằm dưới route
\path|/admin/*|, dùng chung codebase và container với storefront.

\begin{table}[H]
\centering
\caption{Các module chính của Admin Panel}
\label{table:admin-modules}
\begin{compacttable}
\begin{tabularx}{\textwidth}{|L{2.8cm}|L{3.2cm}|Y|}
\hline
\textbf{Module} & \textbf{Route frontend} & \textbf{Endpoint backend chính} \\ \hline
Dashboard & \path|/admin/dashboard| & \path|/api/orders/admin/analytics/*|, \path|/api/inventory/low-stock| \\ \hline
Products & \path|/admin/products| & \path|/api/products| \\ \hline
Inventory & \path|/admin/inventory| & \path|/api/inventory/admin|, \path|/api/inventory/stock-in|, \path|/api/inventory/stock-out| \\ \hline
Orders & \path|/admin/orders| & \path|/api/orders/admin| \\ \hline
Users & \path|/admin/users| & \path|/api/users/admin| \\ \hline
Vouchers & \path|/admin/vouchers| & \path|/api/vouchers| \\ \hline
Flash sales & \path|/admin/flash-sales| & \path|/api/flash-sales| \\ \hline
Content/Reviews & \path|/admin/content/*|, \path|/admin/reviews| & \path|/api/content/*|, \path|/api/reviews/admin| \\ \hline
\end{tabularx}
\end{compacttable}
\end{table}

Các mutation trong Admin Panel được thực hiện bằng Server Actions. Action gọi \texttt{adminFetch}
hoặc \texttt{adminMutate}, phía server đọc access token từ httpOnly cookie và gọi API Gateway.
Sau khi mutation thành công, action \texttt{revalidatePath} để làm mới dữ liệu và redirect về
trang danh sách. Thiết kế này giúp token không xuất hiện trong JavaScript phía trình duyệt,
đồng thời vẫn tận dụng được cơ chế caching/revalidation của Next.js.

% \begin{figure}[H]
%     \centering
%     \includegraphics[width=0.85\textwidth]{DUONG_DAN_ANH}
%     \caption{Storefront: trang chủ, danh sách sản phẩm, giỏ hàng và checkout}
%     \label{fig:storefront-screens}
% \end{figure}

% \begin{figure}[H]
%     \centering
%     \includegraphics[width=0.85\textwidth]{DUONG_DAN_ANH}
%     \caption{Admin Panel: dashboard, products, inventory, orders và flash-sales}
%     \label{fig:admin-panel-screens}
% \end{figure}

Giới hạn của frontend trong phạm vi đồ án được ghi nhận rõ: ảnh sản phẩm/banner nhập bằng URL
thay vì upload file thật; review-service mới hỗ trợ list/delete, chưa có workflow approve/reject
đầy đủ; chức năng ban/unban user chưa tích hợp Keycloak Admin API; dashboard analytics tổng hợp
trực tiếp từ order-service, phù hợp demo nhưng chưa phải kiến trúc analytics quy mô lớn.

% =======================================================================
\section{Docker Compose và vận hành local}
\label{section:5.5}

Môi trường local dùng \path|docker-compose.yml| để khởi động toàn bộ backend và hạ tầng.
Các container được đặt chung network \path|ecommerce-network|; các stateful service dùng
Docker volume để giữ dữ liệu giữa các lần restart. Healthcheck được khai báo cho PostgreSQL,
Redis, Kafka, Elasticsearch, Keycloak, Eureka và các service Spring Boot. Nhờ đó Compose khởi
động theo thứ tự hợp lý thay vì chỉ dựa vào thứ tự khai báo.

% \begin{figure}[H]
%     \centering
%     \includegraphics[width=0.85\textwidth]{DUONG_DAN_ANH}
%     \caption{Deployment diagram Docker Compose local}
%     \label{fig:docker-compose-local}
% \end{figure}

Các endpoint vận hành local quan trọng gồm API Gateway \path|http://localhost:8080|, Eureka
\path|http://localhost:8761|, Keycloak \path|http://localhost:8180|, Prometheus
\path|http://localhost:9090|, Zipkin \path|http://localhost:9411|, Mailpit
\path|http://localhost:8025| và Elasticsearch \path|http://localhost:9200|. Grafana mặc
định dùng cổng \path|3000| trong local backend stack; khi chạy frontend cùng lúc cần remap
hoặc dùng cấu hình production nơi Grafana được đưa sang cổng \path|3001|.

% \begin{figure}[H]
%     \centering
%     \includegraphics[width=0.85\textwidth]{DUONG_DAN_ANH}
%     \caption{Eureka dashboard sau khi các service đăng ký thành công}
%     \label{fig:eureka-dashboard}
% \end{figure}

% =======================================================================
\section{Triển khai production-like trên AWS EC2}
\label{section:5.6}

Môi trường triển khai thực nghiệm sử dụng một máy ảo AWS EC2 single-host tại khu vực
ap-southeast-1. Đây không phải kiến trúc high availability, nhưng phù hợp với phạm vi DATN:
đủ để chứng minh hệ thống có thể build image, chạy container, cấu hình reverse proxy/HTTPS,
giám sát runtime và thực hiện các bài kiểm thử hiệu năng/resilience trên môi trường gần thực tế.

\begin{table}[H]
\centering
\caption{Thành phần triển khai production-like}
\label{table:aws-resources}
\begin{compacttable}
\begin{tabularx}{\textwidth}{|L{3.5cm}|Y|}
\hline
\textbf{Thành phần} & \textbf{Vai trò} \\ \hline
AWS EC2 & Máy chủ single-host chạy Docker Compose production stack \\ \hline
Docker Compose production override & Dùng GHCR image, cấu hình memory limit/reservation, remap Grafana 3001 \\ \hline
Caddy & Reverse proxy, HTTPS tự động qua Let's Encrypt và \texttt{nip.io} \\ \hline
GitHub Actions & Build/push backend và frontend image lên GHCR, hỗ trợ deploy lên EC2 \\ \hline
GHCR & Registry lưu Docker image theo service \\ \hline
\end{tabularx}
\end{compacttable}
\end{table}

File \path|docker-compose.prod.yml| không build image trực tiếp trên EC2 mà tham chiếu image
theo dạng \path|ghcr.io/{owner}/{service}:{tag}|. Cách này giảm thời gian deploy trên
máy chủ, tách build khỏi runtime và giúp cùng một image có thể được kéo lại khi cần rollback.
Production override cũng đặt \path|mem_limit| và \path|mem_reservation| cho từng container,
giúp stack ổn định hơn trong giới hạn tài nguyên của EC2. Caddy định tuyến các subdomain
\path|app.*|, \path|api.*|, \path|keycloak.*| và \path|grafana.*|; Grafana dùng host
port \path|3001| để không xung đột với frontend.

% \begin{figure}[H]
%     \centering
%     \includegraphics[width=0.85\textwidth]{DUONG_DAN_ANH}
%     \caption{Deployment diagram AWS EC2 production-like stack}
%     \label{fig:aws-deployment}
% \end{figure}

% \begin{figure}[H]
%     \centering
%     \includegraphics[width=0.85\textwidth]{DUONG_DAN_ANH}
%     \caption{GitHub Actions build/push image lên GHCR thành công}
%     \label{fig:github-actions}
% \end{figure}

% =======================================================================
\section{Chiến lược kiểm thử}
\label{section:5.7}

Kiểm thử được tổ chức theo nhiều tầng. Tầng thấp nhất là unit test và integration test cho
logic nghiệp vụ quan trọng. Tầng tiếp theo là backend readiness smoke test, chạy theo các
suite tách nhỏ để tránh quá tải tài nguyên và giúp khoanh vùng lỗi nhanh. Tầng cao hơn là
kiểm thử hiệu năng bằng k6 trên môi trường AWS EC2 và kiểm thử resilience bằng cách chủ động
dừng service/hạ tầng rồi quan sát khả năng tự phục hồi.

\begin{table}[H]
\centering
\caption{Kết quả backend readiness smoke test}
\label{table:backend-readiness}
\begin{compacttable}
\begin{tabularx}{\textwidth}{|L{3cm}|L{3.1cm}|Y|}
\hline
\textbf{Suite} & \textbf{Kết quả} & \textbf{Artifact} \\ \hline
Maven & PASS, 5 pass, 0 fail & \path|01-maven-PASS.md| \\ \hline
Infra & PASS, 27 pass, 0 fail & \path|02-infra-PASS.md| \\ \hline
Auth/user & PASS, 30 pass, 0 fail & \path|03-auth-user-PASS-after-fix.md| \\ \hline
Catalog/search & PASS, 49 pass, 0 fail & \path|04-catalog-search-PASS.md| \\ \hline
Cart/voucher & PASS, 57 pass, 0 fail & \path|05-cart-voucher-PASS-after-fix.md| \\ \hline
Order COD/review & PASS, 67 pass, 0 fail & \path|06-order-cod-review-PASS.md| \\ \hline
VNPAY & PASS, 55 pass, 0 fail & \path|07-vnpay-PASS.md| \\ \hline
Flash-sale & PASS, 62 pass, 0 fail & \path|08-flash-sale-PASS-after-fix.md| \\ \hline
Security & PASS, 76 pass, 0 fail, 1 warn & \path|09-security-PASS-trimmed-memory.md| \\ \hline
\end{tabularx}
\end{compacttable}
\end{table}

Kết quả tổng hợp cho thấy 9/9 split suites đạt. Các lỗi phát hiện trong quá trình chuẩn bị
kiểm thử, như healthcheck bị 401, guest cart session id bị reject, lỗi serializer cache sản
phẩm và cấu hình Metaspace cho flash-sale test, đã được sửa trước khi ghi nhận kết quả final.

% =======================================================================
\section{Kết quả kiểm thử hiệu năng}
\label{section:5.8}

Các bài kiểm thử hiệu năng được thực hiện trên môi trường AWS EC2 chạy Docker Compose
production-like stack. Load generator là máy cá nhân chạy k6. Dataset thực nghiệm gồm 100 sản
phẩm, 500 user/token và một flash-sale campaign có 100 slot. Nguyên tắc đưa số liệu vào báo
cáo là mọi metric phải đối chiếu được với raw artifact trong thư mục \texttt{.test/results}.

\begin{table}[H]
\centering
\caption{Tổng hợp kết quả kiểm thử hiệu năng bằng k6}
\label{table:perf-results}
\begin{compacttable}
\begin{tabularx}{\textwidth}{|L{3.1cm}|L{2cm}|L{1.5cm}|L{2.2cm}|L{2.2cm}|Y|}
\hline
\textbf{Kịch bản} & \textbf{Thời lượng} & \textbf{Max VU} & \textbf{Requests} & \textbf{p95} & \textbf{Lỗi} \\ \hline
Catalog soak & 34m03s & 200 & 220,956 & 61.68 ms & 0.00\% \\ \hline
Checkout stress & 7m00s & 50 & 13,734 & 160.92 ms & 0.00\% \\ \hline
Flash-sale spike & 1m30s & 500 & 78,448 & 1.04 s overall & 0 duplicate buyer \\ \hline
\end{tabularx}
\end{compacttable}
\end{table}

\textbf{Catalog soak} mô phỏng người dùng duyệt danh sách sản phẩm trong 34 phút 03 giây với
tải tối đa 200 VU. Kịch bản tạo 55,239 iterations và 220,956 HTTP requests, p95 đạt 61.68 ms,
p99 đạt 85.48 ms, error rate 0.00\%. Kết quả này cho thấy read path catalog/search/cache đáp
ứng tốt trên single-host trong phạm vi dataset đồ án.

\textbf{Checkout stress} mô phỏng tạo đơn hàng với tải tối đa 50 VU trong 7 phút. Kịch bản
ghi nhận 6,817 iterations, 13,734 HTTP requests, order success check 100.00\%, HTTP failure
rate 0.00\% và p95 latency 160.92 ms. Đây là luồng nặng hơn catalog vì đi qua product,
inventory, voucher, order, Kafka và notification, nhưng vẫn duy trì độ trễ phù hợp với mục
tiêu thực nghiệm.

\textbf{Flash-sale spike} mô phỏng 500 VU trong 90 giây tranh mua 100 slot. Kết quả có đúng
100 purchase thành công trên 100 slot, \texttt{soldCount} của campaign bằng 100, số đơn
confirmed cho campaign bằng 100 và không có duplicate buyer. Bài test ghi nhận 36,134 phản
hồi sold-out; p95 tổng thể là 1.04 giây, còn p95 của nhóm expected responses là 335.55 ms.
Điều này chứng minh trong kịch bản thực nghiệm đã chạy, Redis Lua atomic reservation kết hợp
buyer set, compensation và reconciliation không gây over-sell.

% \begin{figure}[H]
%     \centering
%     \includegraphics[width=0.85\textwidth]{DUONG_DAN_ANH}
%     \caption{Kết quả k6 catalog, checkout và flash-sale spike}
%     \label{fig:perf-results}
% \end{figure}

% =======================================================================
\section{Kết quả kiểm thử khả năng chịu lỗi}
\label{section:5.9}

Để đánh giá resilience, bốn kịch bản được thực hiện bằng cách chủ động gây lỗi vào service
hoặc hạ tầng. Mục tiêu không phải chứng minh hệ thống có high availability hoàn chỉnh, mà là
kiểm chứng các cơ chế đã thiết kế: circuit breaker/fallback, outbox replay, graceful
degradation và Saga compensation.

\begin{table}[H]
\centering
\caption{Tổng hợp kết quả kiểm thử khả năng chịu lỗi}
\label{table:chaos-results}
\begin{compacttable}
\begin{tabularx}{\textwidth}{|L{3.2cm}|Y|L{2.1cm}|}
\hline
\textbf{Kịch bản} & \textbf{Kết quả chính} & \textbf{Trạng thái} \\ \hline
Kill order-service & Trong 30 giây downtime, order-created đạt 47.15\%, HTTP failure 25.94\%; sau restart gateway/order health trả 200 & Completed \\ \hline
Kill Kafka & Khi Kafka down, order vẫn được tạo, outbox row ở trạng thái pending; sau Kafka restart, event được process và order kết thúc \texttt{CONFIRMED} & PASS \\ \hline
Kill Redis & Product list vẫn trả 200 trong Redis outage; cart phụ thuộc Redis degrade 500 trong outage và recover 200 sau restart & PASS \\ \hline
Inventory failure compensation & Order \path|4170422f-fc2d-49f8-aa5f-c7cd1d79266a| chuyển \path|CANCELLED|, inventory còn \path|quantity=0|, \path|reserved_quantity=0| & PASS \\ \hline
\end{tabularx}
\end{compacttable}
\end{table}

Kịch bản Kill Kafka là bằng chứng trực tiếp cho Transactional Outbox: service nghiệp vụ vẫn
ghi được dữ liệu và ghi event pending vào database khi broker không khả dụng; sau khi broker
khởi động lại, poller phát lại event và Saga tiếp tục đến trạng thái cuối. Kịch bản Kill Redis
cho thấy hệ thống có phân biệt read path có fallback (product list) với chức năng phụ thuộc
trực tiếp Redis (cart/rate-limit/flash-sale). Cách ghi nhận này trung thực với bản chất hệ
thống: một số luồng có graceful degradation, một số luồng phải fail fast để tránh sai dữ liệu.

% \begin{figure}[H]
%     \centering
%     \includegraphics[width=0.85\textwidth]{DUONG_DAN_ANH}
%     \caption{Bằng chứng resilience: Kafka outbox replay và Redis recovery}
%     \label{fig:resilience-evidence}
% \end{figure}

% =======================================================================
\section{Đánh giá tổng thể}
\label{section:5.10}

\subsection{Đối chiếu với mục tiêu đề tài}

\begin{table}[H]
\centering
\caption{Đối chiếu kết quả với mục tiêu ban đầu}
\label{table:objectives-eval}
\begin{compacttable}
\begin{tabularx}{\textwidth}{|L{3.5cm}|Y|L{2cm}|}
\hline
\textbf{Mục tiêu} & \textbf{Kết quả đạt được} & \textbf{Đánh giá} \\ \hline
Phân rã hệ thống theo microservices & 13 business service, 3 Spring Cloud service, database/cache/search tách theo ownership & Đạt \\ \hline
Áp dụng distributed patterns & Saga Orchestration, Transactional Outbox, Idempotent Consumer, CQRS-lite, scheduler/reconciliation & Đạt \\ \hline
Bảo mật tập trung & Keycloak, JWT validation tại Gateway, trusted headers, RBAC, rate limiter & Đạt \\ \hline
Thanh toán trực tuyến & VNPAY sandbox, HMAC-SHA512, Return URL, IPN, timeout scheduler & Đạt \\ \hline
Flash sale chống over-sell & Redis Lua, buyer set, compensation, reconciliation; spike test 500 VU xác nhận 100/100 slot và 0 duplicate buyer & Đạt \\ \hline
Triển khai và quan sát & Docker Compose local/AWS EC2, Caddy HTTPS, GitHub Actions/GHCR, Prometheus/Grafana/Zipkin & Đạt trong phạm vi DATN \\ \hline
Giao diện demo hệ thống hoàn chỉnh & Next.js storefront và Admin Panel với các module quản trị chính & Đạt mức demo DATN \\ \hline
\end{tabularx}
\end{compacttable}
\end{table}

\subsection{Hạn chế}

Hệ thống đã đạt mục tiêu của đồ án nhưng vẫn có các giới hạn cần ghi nhận. Thứ nhất, môi
trường triển khai là single-host EC2 nên chưa có high availability; nếu EC2 gặp sự cố, toàn
bộ hệ thống bị ảnh hưởng. Thứ hai, CI/CD hiện đáp ứng build/push/deploy cơ bản, chưa có
canary/blue-green deployment, security scan bắt buộc hay quản lý secret chuyên dụng. Thứ ba,
frontend/Admin Panel phục vụ tốt demo DATN nhưng chưa phải sản phẩm thương mại hoàn chỉnh:
chưa có upload file qua object storage, workflow kiểm duyệt review đầy đủ, ban/unban user qua
Keycloak và analytics service riêng. Thứ tư, event schema đang dùng module \texttt{common},
phù hợp monorepo đồ án nhưng production lớn nên cân nhắc schema registry hoặc contract testing.
Thứ năm, hệ thống có metrics/tracing nhưng chưa có log aggregation và alerting production-grade.

\end{document}
